\documentclass[UTF8, 11pt, a4paper]{ctexart}
% 页面边距设置
\usepackage[left=2.5cm, right=2.5cm, top=3cm, bottom=3cm]{geometry}
\usepackage{xcolor, fancyhdr, titlesec, hyperref, amsmath, amssymb, amsthm}
\usepackage[many]{tcolorbox}
\usepackage{enumitem} % 用于自定义列表间距

% ==========================================
% 颜色定义 (绝对没有绿色！采用牛津蓝与勃艮第红的经典搭配)
% ==========================================
\definecolor{OxfordBlue}{RGB}{0, 33, 71}      % 主色调：深邃的牛津蓝
\definecolor{Burgundy}{RGB}{128, 0, 32}       % 强调色：勃艮第酒红 (用于定理、警告、公式)
\definecolor{LightGrey}{RGB}{248, 249, 250}   % 框体背景色：极浅的灰色
\definecolor{Highlight}{RGB}{255, 243, 205}   % 行内高亮底色：柔和的淡黄色

% ==========================================
% 自定义行内文本命令
% ==========================================
\newcommand{\hl}[1]{\colorbox{Highlight}{#1}} % 行内高亮文字
\newcommand{\imp}[1]{\textbf{\color{Burgundy}#1}} % 重点强调文字(红色加粗)

% --- 超链接设置 ---
\hypersetup{colorlinks=true, linkcolor=OxfordBlue, urlcolor=OxfordBlue}

% ==========================================
% 页眉页脚设置
% ==========================================
\pagestyle{fancy}
\fancyhf{}
\fancyhead[L]{\color{OxfordBlue}\textbf{《网页深度阅读笔记》}}
\fancyhead[R]{\color{OxfordBlue}\leftmark}
\fancyfoot[C]{\thepage}
\renewcommand{\headrulewidth}{0.8pt}
\renewcommand{\headrule}{\hbox to\headwidth{\color{OxfordBlue}\leaders\hrule height \headrulewidth\hfill}}

% ==========================================
% 标题美化
% ==========================================
\titleformat{\section}{\Large\bfseries\color{OxfordBlue}}{\thesection}{1em}{}[\titlerule]
\titleformat{\subsection}{\large\bfseries\color{OxfordBlue!85!black}}{\thesubsection}{1em}{}

% ==========================================
% 各类特色笔记框定义
% ==========================================

% 1. 定义/概念框 (蓝色基调，带阴影)
\newtcolorbox{conceptbox}[1][]{
    enhanced, colback=LightGrey, colframe=OxfordBlue, fonttitle=\bfseries,
    title=#1, attach boxed title to top left={yshift=-2mm, xshift=5mm},
    boxed title style={colback=OxfordBlue, sharp corners=all, rounded corners=northeast},
    arc=2mm, drop fuzzy shadow
}

% 2. 定理/公式框 (红色基调，突出视觉重点)
\newtcolorbox{theorembox}[1][]{
    enhanced, colback=Burgundy!5, colframe=Burgundy, fonttitle=\bfseries,
    title=#1, attach boxed title to top left={yshift=-2mm, xshift=5mm},
    boxed title style={colback=Burgundy, sharp corners=all, rounded corners=northeast},
    arc=2mm, drop fuzzy shadow
}

% 3. 原文摘抄框 (左侧粗线，极简风格)
\newtcolorbox{quotebox}{
    colback=OxfordBlue!3, colframe=OxfordBlue,
    leftrule=4pt, rightrule=0pt, toprule=0pt, bottomrule=0pt, arc=0pt
}

% 4. 注意事项/警告框 (虚线边框，红色)
\newtcolorbox{warningbox}{
    enhanced, colback=white, colframe=Burgundy,
    boxrule=1pt, arc=1mm, borderline={0.5pt}{2pt}{Burgundy, dashed},
    title=\textcolor{Burgundy}{\textbf{⚠️ 易错点提示}}, coltitle=Burgundy,
    attach title to upper, after title={\par\smallskip}
}

% ==========================================
% 正文开始
% ==========================================
\title{\vspace{-2cm}\color{OxfordBlue}\Huge\textbf{网页文献阅读与重构笔记}}
\author{Autumn}
\date{\today}

\begin{document}
\maketitle
\tableofcontents
\thispagestyle{empty}
\newpage

\section{第一章：信息论与贝叶斯推断}
我们在阅读这篇网页时，作者花了大篇幅讨论概率论的基础。本章主要梳理其中的数学逻辑和关键论点。

% --- 概念框示例 ---
\begin{conceptbox}[核心定义：条件概率 (Conditional Probability)]
条件概率是指事件 A 在另外一个事件 B 已经发生条件下的发生概率。表示为 $P(A|B)$，读作“在 B 发生的条件下 A 发生的概率”。
\end{conceptbox}

在理解了上述定义后，网页作者引用了拉普拉斯的一句名言，我觉得总结得非常好：

% --- 摘抄框示例 ---
\begin{quotebox}
“概率论只不过是把常识用数学公式严谨地表达了出来。它让我们在面对不确定性时，依然能够做出最优的决策。” —— \textit{网页第2自然段}
\end{quotebox}

\subsection{1.1 核心公式定理}
接下来是全篇最核心的部分。作者不仅给出了公式，还给出了推导。为了在复习时一眼看到重点，我把它放进了独立的公式定理框中：

% --- 定理/公式框示例 ---
\begin{theorembox}[定理：贝叶斯公式 (Bayes' Theorem)]
假设 $A$ 和 $B$ 是两个随机事件，且 $P(B) \neq 0$，则有：
\begin{equation}
    P(A|B) = \frac{P(B|A) P(A)}{P(B)}
\end{equation}
其中，$P(A)$ 称为\imp{先验概率} (Prior)，$P(A|B)$ 称为\imp{后验概率} (Posterior)。
\end{theorembox}

% --- 多行公式推导示例 ---
\subsection{1.2 公式的推导过程}
根据全概率公式，我们可以对分母进行展开推导。以下是详细的步骤（注意对齐）：
\begin{align}
    P(A|B) &= \frac{P(A \cap B)}{P(B)} \notag \\ % notag 表示这一行不编号
           &= \frac{P(B|A) P(A)}{\sum_{i=1}^{n} P(B|A_i)P(A_i)} \\
           &\propto P(B|A) P(A)
\end{align}
从式 (3) 可以看出，\hl{后验概率正比于似然函数乘以先验概率}。这句话是整篇文章的灵魂。

% --- 列表与警告框示例 ---
\subsection{1.3 实际应用步骤}
在将贝叶斯推断应用到实际代码中时，作者给出了以下几个步骤：
\begin{itemize}[itemsep=3pt, parsep=0pt]
    \item \textbf{Step 1}: 确定参数的先验分布模型。
    \item \textbf{Step 2}: 根据观测数据写出似然函数 (Likelihood)。
    \item \textbf{Step 3}: 利用积分或马尔可夫链蒙特卡洛 (MCMC) 方法求解后验分布。
\end{itemize}

\begin{warningbox}
在执行 Step 3 时，分母的积分往往是无法解析求解的（Intractable）。千万不要尝试强行手算高维分布的连续积分，直接使用 \texttt{PyMC3} 或 \texttt{Stan} 等概率编程库进行近似采样。
\end{warningbox}

\end{document}
